% =====================================================================================================
% Online Appendix E. Observational production functions and algorithmic progress.
% v2 (writer appendix_E, 2026-09-24). Replaces v1 Section VI and the observational/progress parts of v1 Appendix D.
% Sources: output/memos/ra4_obsfix(.md, _review.md) [binding where it differs from m4/m5]; m4_observational, m5_progress,
% m6_montecarlo memos and reviews; theory from paper/sections/appendix_proofs.tex (Props. A3-A7, Cors. A2, A4, A5).
% Numbers verified against output/tables/ra4_obsfix_*.csv, m4_observational_*.csv, m5_progress_*.csv,
% m6_montecarlo_industry_summary.csv and data/processed/ra4_obsfix/*.json (see paper/notes/writer2_appendix_E.md).
% Figures: ra4_obsfix_fig_matched, ra4_obsfix_fig_allocative, appE_ridge (code/paper/make_appendix_e_ridge.py; m5 ridge relabeled).
% Tables: paper/tables/appE_*.tex (copies of output/tables/ra4_obsfix_tabE*.tex and v1 appD_attenuation/ceg_sahal,
% with notation and cross-references updated for v2).
% =====================================================================================================

\section{Observational Production Functions and Algorithmic Progress}\label{app:obs}

The main text reads developers' choices through a technology measured in designed experiments. This appendix asks the
reverse question: what do released models, used as production data, reveal about the technology and about its change
over time? Much of what is known about returns to compute at scale, and about the pace of algorithmic progress, comes
from such data: regressions of benchmark scores on the inputs of models released by many developers
\citep{ruan2024observational,maiapolo2024sloth,owen2024predictable,mertens2026secret}, and of perplexity on the inputs
and publication dates of published models \citep{ho2024algorithmic}. In these data the firms chose the inputs, and
\citet{whitfill2025note} and \citet{konig2026validity} warn that the estimates may be confounded.
Section~\ref{app:obs-theory} collects the results of Online Appendix~\ref{app:proofs} that apply to such data, and the
remaining sections take them to the data.

There are four findings. First, once the design is matched, cross-lab returns to compute on HellaSwag do not differ
detectably from the experimental technology. This is a failure to reject in a low-power joint test, not evidence of
unbiasedness, and the split between parameters and data is not reliably estimated. Second, compute is measured well
enough that errors in variables do not matter, but the standard remedies for transmission bias (instruments, dynamic
panels, proxy inversion) cannot be applied to these data. Third, in the data of \citet{ho2024algorithmic} the rate of
algorithmic progress is weakly identified: estimates lie on a ridge along which the bias of technical change trades off
against its rate, and the published point is not the minimizer of the authors' objective. Fourth, counting only the
elimination of under-training, the move from Kaplan-rule to Chinchilla-rule allocations raised the training-cost
efficiency of large models by a factor of 2.0 under the reference technology of Section~\ref{sec:wedge} (1.1 to 4.9
across technologies). The tenfold gain of \citet{gundlach2025origin} is the counterfactual cost of keeping the earlier
rule at 2025 frontier compute, not a gain realized between eras.

% -----------------------------------------------------------------------------------------------------
\subsection{What Cross-Lab Data Can Identify}\label{app:obs-theory}

The results below, proved in Online Appendix~\ref{app:proofs}, concern a common technology with Hicks-neutral
productivity $\omega$ that the developer knows when it chooses inputs, and log output $y=\omega+F(n,d)-\epsilon$, where
$F$ is the technology in output units and $\epsilon$ is noise.

\emph{Transmission.}---If compute responds to productivity, $c=\pi_0+\pi_1\omega+\eta$, the least-squares slope of output
on compute is biased with the sign of $\pi_1$ (Proposition~\ref{prop:A-transmission}). The bias is upward when more
productive developers raise more compute (funding, $\pi_1>0$) and downward when every developer trains to a common
capability target ($\pi_1=-1/\gamma$); if the target is hit exactly, cross-lab data carry no information about the
frontier elasticity $\gamma$. The forward and reverse regressions bracket $\gamma$ under target and exogenous budgets but
not under funding, so the errors-in-variables bracket of \citet{klepper1984consistent} does not survive transmission
from productivity to compute. Hicks-neutral transmission leaves the allocation exponent $a$ unbiased, because
productivity does not move the input mix at given compute, and it compresses measured productivity dispersion.

Proposition~\ref{prop:A-transmission} places runs on the expansion path. Released models are far off it
(Section~\ref{sec:wedge}), and their log wedges enter output through the Farrell loss of Lemma~\ref{lem:ce}. Two
consequences follow by linear projection. When the regression is run on both inputs, the Farrell loss is part of the
technology $F$, and transmission loads on the mix only if developers' wedges covary with their productivity at given
compute. If they do not, the projection of $\omega$ on $(n,d)$ runs along $c=\ln6+n+d$, and transmission shifts the
two input elasticities by the same amount. If more productive developers also over-train more (or less), the two
biases differ and can have opposite signs. When the regression is run on compute alone, it imposes equal weights on
the two inputs, and the productivity residual inherits the Hall-type term of Corollary~\ref{cor:hall}
(Section~\ref{app:obs-matched}).

\emph{Selection.}---If a model is released only when its output clears a threshold, least squares on released models is
attenuated. This is the classical result for truncation on the dependent variable
\citep{goldberger1981linear,hausman1977social}. What Proposition~\ref{prop:selection} adds is the interaction with
transmission. Under funding, mild selection leaves the upward transmission bias dominant and severe selection reverses
it. Under target or exogenous budgets the net bias is negative, and the forward--reverse bracket survives in released
samples.

\emph{Proxies.}---Proxy estimators \citep{olley1996dynamics,levinsohn2003estimating,ackerberg2015identification} recover
productivity by inverting a policy function, so they need a state variable, a flexible input and a timing assumption.
Here the natural state variable is the developer's compute capacity, the provisioned cluster, which is fixed before a
run. The flexible input is the run's split of its compute budget between $N$ and $D$, chosen at the start of the run
given capacity and the recipe's quality. Release follows the observed outcome. The flexible input cannot proxy
Hicks-neutral productivity, because at given compute the policy function is flat in $\omega$
(Proposition~\ref{prop:proxy}(ii)). Compute itself fails as a Levinsohn--Petrin-type proxy for a different reason.
Under a strictly monotone budget rule the inputs on the path are functions of $c$ alone, so a control function in $c$
absorbs the technology: this is the functional-dependence problem of \citet{ackerberg2015identification}
(Proposition~\ref{prop:proxy}(iii)). Identification would then require timing, for example compute committed before a
recipe innovation, or an observed capacity investment, which public data do not record. Under a capability target the
mix chosen for a given loss does co-move with $\omega$ when $\alpha\neq\beta$, but only through compute
(Lemma~\ref{lem:alloc}(v)). Conditional on $c$, the mix remains uninformative about $\omega$.

\emph{Technical change.}---\citet{diamond1978measurement} show that time series of outputs, inputs and marginal
products cannot separate the elasticity of substitution from the bias of technical change. Time series of
compute-optimal models still identify the sign of the bias, though not its size. Because the FLOP accounting fixes the relative cost of a
parameter and a token, the compute-optimal ratio $M^*(C,t)$ drifts at given compute as
$\partial\ln M^*/\partial t=-\mathcal B_t\,\sigma^*_t/(1-\sigma^*_t)$. Here $\mathcal B_t$, the Hicks bias, is the
change over time in the marginal rate of technical substitution at fixed inputs (Proposition~\ref{prop:A-dmr}(i)).
Since $0<\sigma^*<1$ at interior optima (Lemma~\ref{lemma:sigma}), $M^*$ falls if and only if technical change is
parameter-using ($\mathcal B_t>0$). In the generalized family, with inner exponents $a_1$ and $b_1$ and factor augmentation at rates
$g_N$ and $g_D$, the bias is $\mathcal B=b_1g_D-a_1g_N$, which is $\beta g_D-\alpha g_N$ when $\kappa=1$. The drift
therefore signs $\beta g_D-\alpha g_N$, not $g_D-g_N$. For example, with $(\alpha,\beta)=(0.30,0.45)$ and
$(g_N,g_D)=(0.5,0.4)$, $\mathcal B=0.03>0$ and the compute-optimal $D/N$ falls, although parameters are augmented faster.
Within the family $g_N$ and $g_D$ are identified separately, but only through its functional form; the curvature, and
with it the size of the bias, is not. The result is methodological. We do not take it to published compute-optimal laws
of different vintages, which differ in tokenizers, parameter-counting conventions and outputs. The sign result also
requires that the effective relative cost of tokens not drift through data costs, data constraints or serving demand
(Proposition~\ref{prop:wedge}). A designed IsoFLOP sweep at each vintage on a common output would identify the sign,
and through Proposition~\ref{prop:modelfree} the size.

\emph{Allocative versus technical change.}---A model's log excess compute over best practice splits into an allocative
term $\Phi(w)=\ln(C/C_{\min})$, which depends only on its wedge (Lemma~\ref{lem:ce}), and a technical term, the
productivity gap in compute-equivalent units (Proposition~\ref{prop:farrell}). A reallocation at fixed productivity,
such as the move from Kaplan-rule to Chinchilla-rule allocations, is purely allocative. Three corollaries concern
growth accounting. First, the compute-equivalent gain of \citet{gundlach2025origin} between two technologies is the
same at every compute level if and only if they share the irreducible loss $E$ and the frontier elasticity $\gamma$
(Corollary~\ref{cor:ceg}). Second, regressions of frontier loss on compute that omit time effects attribute part of
algorithmic progress to scale, overstating the compute elasticity by the factor $1/(1-s_A)$, where $s_A$ is the
algorithmic share of effective-compute growth (Corollary~\ref{cor:sahal}; \citealp{nagy2013statistical}). Third, a
productivity residual computed with equal weights on $\ln N$ and $\ln D$, as $C=6ND$ suggests, misattributes
$\tfrac12(\varepsilon_N-\varepsilon_D)(\Delta n-\Delta d)$ to productivity. When over-trained models ($w>1$) shift
further toward tokens, the residual understates progress (Corollary~\ref{cor:hall}; \citealp{hall1988relation}).

% -----------------------------------------------------------------------------------------------------
\subsection{Returns to Compute: A Design-Matched Validation}\label{app:obs-matched}

\emph{Data and specification.}---The observational sample has 128 dense base models from 38 families and 22 developers,
released between March 2021 and July 2024 (Online Appendix~\ref{app:data}). The models share no validation loss, so
output is the log-odds of above-chance accuracy, $q=\ln[\tilde a/(1-\tilde a)]$, where $\tilde a$ is accuracy net of
chance, rescaled to the unit interval and clipped to $[0.01,0.99]$. The main benchmark is HellaSwag
\citep{zellers2019hellaswag}, evaluated 10-shot on the Open LLM Leaderboard. We estimate
$q_i=\theta_0+\theta_Nn_i+\theta_Dd_i+\omega_i+e_i$, which is Cobb--Douglas in the odds with Hicks-neutral productivity
$\omega_i$, and the regression of $q_i$ on $c_i$ alone. Its slope $\theta_C$ is what observational scaling studies
report.

\emph{Why the design must be matched.}---A comparison with global experimental elasticities misleads, because the
technology is curved and released models sit in a different part of it. In the designed sweeps, HellaSwag log-odds are
strongly concave in $(n,d)$; the $F$-statistic for the quadratic terms is 70 in the sweep of \citet{gadre2024language}.
On the experimental surface described below, at $N=1$B the local data elasticity falls from 0.57 at 20 tokens per
parameter to 0.07 at 500. Released models are heavily over-trained (median $D/N$ of 213, ninetieth percentile 1,715), so
they sit where an extra token is worth little. Global log-linear fits give observational data elasticities of 0.19 to
0.26, against 0.44 in the OLMo ladder and 0.51 in \citet{gadre2024language}, but much of that gap is the technology,
not bias.

We therefore match the design. We fit a quadratic surface $\hat f(n,d)$ with design intercepts to the designed runs of
the OLMo ladder \citep{bhagia2024establishing} and of \citet{gadre2024language} with $N\geq0.1$B, 86 runs in all
($R^2=0.994$). We keep the released models inside the convex hull of the experimental inputs, the common-support
trimming of \citet{dehejia1999causal}. We then apply each estimator twice, with the same fixed effects and controls: to
the observed outputs, and to the outputs $q^*_i=\hat f(n_i,d_i)$ that the experimental technology produces with the
same inputs. The difference is the estimator's bias. The benchmark comes from other developers' recipes, not from the
treated population. The exercise is therefore a validation across populations, in the sense of
\citet{todd2006assessing} and \citet{hotz2005predicting}, rather than a comparison with an experiment on the same units
as in \citet{lalonde1986evaluating}. The null hypothesis is joint: a common technology up to Hicks-neutral shifts,
outputs that map one to one across evaluation harnesses, and no transmission or selection bias. Section~\ref{sec:tech} documents factor-biased differences across
recipes, which would violate the first part, so a null bias can also reflect offsetting failures.

\emph{Support.}---An earlier version of this paper included the released OLMo-2 7B and 13B models
\citep{teamolmo2024olmo} among the benchmark runs. Their developer chose their inputs, so they are no longer part of the
experimental benchmark. The cost is support. Without them the largest designed $D$ is 921B tokens (one run of
\citealp{gadre2024language} at 1.4B parameters), and above 3.8B parameters it is 634B. The strict support contains
20 released models from 9 families and 6 developers, all released between March 2021 and September 2023 with
$N\le6.7$B and $D\le627$B; five of the nine families have two or more models, 16 models in all. The benchmark that
includes OLMo-2 (88 runs) covers 57 models from 30 families and 19 developers, with $N\le8.8$B. Fifteen of its families
have two or more models (42 models) and fifteen are singletons. We report both designs.

\emph{Inference.}---Inference must allow for few clusters and for the benchmark's own uncertainty. Standard errors are
CV3 jackknife standard errors by developer \citep{mackinnon2023cluster}. The benchmark's uncertainty comes from a
design-conditional wild bootstrap of the surface, clustered by design and model size. For each bias we compute three
few-cluster $p$-values that include the benchmark's uncertainty: CV3 with $t(G-1)$ critical values, and restricted
wild-cluster bootstrap-$t$ tests studentized by CR1 or by CV3
\citep{cameron2008bootstrap,webb2023reworking}. Table~\ref{tab:app-matched} reports the largest of the three, so a
rejection requires all three to reject. The effective number of clusters \citep{carter2017asymptotic} is 2.5 to 6.7 of
19 on the 57-model design. On the strict design it is 1.2 to 3.8 of 6; there the three procedures disagree by up to a
factor of ten, and none is reliable.

\begin{table}[tp]
\centering
\caption{Design-Matched Comparison of Observational and Experimental Elasticities}
\label{tab:app-matched}
\footnotesize
\setlength{\tabcolsep}{4pt}
\begin{tabular}{@{}l ccc c c@{}}
\toprule
 & \multicolumn{3}{c}{Compute only, $\theta_C$} & $\theta_N$ & $\theta_D$ \\
\cmidrule(lr){2-4}\cmidrule(lr){5-5}\cmidrule(lr){6-6}
 & Obs. & Exp. & Bias & Bias & Bias \\
 & (1) & (2) & (3) & (4) & (5) \\
\midrule
\multicolumn{6}{@{}l}{\textit{Panel A. HellaSwag, strict benchmark (20 models, 9 families, 6 developers)}} \\
OLS & 0.406 & 0.430 & $-$0.024 & 0.105 & $-$0.090 \\
 & (0.043) & (0.007) & (0.076) [0.77] & (0.107) [0.37] & (0.619) [0.89] \\
Year FE & 0.460 & 0.437 & 0.024 & 0.022 & 0.025 \\
 & (0.035) & (0.007) & (0.080) [0.78] & (0.147) [0.89] & (0.727) [0.97] \\
Developer FE & 0.449 & 0.436 & 0.013 & 0.066 & $-$0.030 \\
 & (0.014) & (0.008) & (0.044) [0.77] & (0.148) [0.68] & (0.763) [0.97] \\
Family FE & 0.473 & 0.464 & 0.008 & 0.185 & $-$0.227 \\
 & (0.243) & (0.010) & (0.158) [0.96] & (0.046) [0.14] & (0.116) [0.15] \\
Lab $\times$ period FE & 0.476 & 0.440 & 0.037 & 0.053 & 0.023 \\
 & (0.239) & (0.008) & (0.132) [0.79] & (0.149) [0.74] & (0.339) [0.95] \\
\addlinespace[3pt]
\multicolumn{6}{@{}l}{\textit{Panel B. HellaSwag, benchmark with OLMo-2 (57 models, 30 families, 19 developers)}} \\
OLS & 0.328 & 0.329 & $-$0.001 & 0.079 & $-$0.050 \\
 & (0.046) & (0.007) & (0.041) [0.98] & (0.091) [0.43] & (0.070) [0.58] \\
Year FE & 0.319 & 0.348 & $-$0.029 & 0.101 & $-$0.113 \\
 & (0.076) & (0.007) & (0.037) [0.47] & (0.107) [0.36] & (0.212) [0.60] \\
Developer FE & 0.393 & 0.370 & 0.024 & 0.141 & $-$0.090 \\
 & (0.054) & (0.008) & (0.018) [0.66] & (0.070) [0.14] & (0.098) [0.57] \\
Family FE & 0.441 & 0.409 & 0.031 & 0.203 & $-$0.262 \\
 & (0.048) & (0.013) & (0.064) [0.63] & (0.059) [0.048] & (0.106) [0.027] \\
Lab $\times$ period FE & 0.442 & 0.405 & 0.037 & 0.165 & $-$0.113 \\
 & (0.032) & (0.010) & (0.039) [0.36] & (0.077) [0.13] & (0.241) [0.65] \\
Notability-propensity control & 0.335 & 0.350 & $-$0.015 & 0.076 & $-$0.057 \\
 & (0.071) & (0.006) & (0.043) [0.80] & (0.125) [0.55] & (0.072) [0.53] \\
\addlinespace[3pt]
\multicolumn{6}{@{}l}{\textit{Panel C. OLS, other outputs and supports}} \\
HellaSwag, strict surface, 57 & 0.328 & 0.314 & 0.014 & 0.090 & $-$0.032 \\
 & (0.046) & (0.014) & (0.045) [0.76] & (0.091) [0.37] & (0.073) [0.73] \\
ARC-Challenge, 20 & 0.466 & 0.473 & $-$0.007 & 0.218 & $-$0.123 \\
 & (0.117) & (0.018) & (0.248) [0.98] & (0.188) [0.50] & (1.429) [0.93] \\
ARC-Challenge, 57, with OLMo-2 & 0.449 & 0.413 & 0.036 & 0.171 & $-$0.045 \\
 & (0.055) & (0.027) & (0.069) [0.61] & (0.182) [0.49] & (0.129) [0.74] \\
Winogrande, 20 & 0.371 & 0.550 & $-$0.179 & $-$0.040 & $-$0.251 \\
 & (0.055) & (0.038) & (0.263) [0.53] & (0.165) [0.90] & (1.341) [0.86] \\
Winogrande, 57, with OLMo-2 & 0.351 & 0.501 & $-$0.150 & $-$0.027 & $-$0.225 \\
 & (0.047) & (0.044) & (0.077) [0.084] & (0.194) [0.89] & (0.202) [0.28] \\
\bottomrule
\end{tabular}

\begin{tablenotes}
Output: log-odds of above-chance accuracy. Obs.: estimator on the observed outputs of released base models. Exp.: the same estimator, with the same fixed effects and controls, on the outputs that the experimental technology produces for the same models: a quadratic in $(\ln N,\ln D)$ fitted to the OLMo ladder and \citet{gadre2024language} ($N\geq0.1$B; 86 runs); Panel B adds the released OLMo-2 7B and 13B models (88 runs), as an earlier version of this paper did. Each panel uses the models inside the convex hull of its own runs; `strict surface, 57' evaluates the Panel A surface on the Panel B models, 37 of which lie outside its support. Bias $=$ Obs.\ $-$ Exp. Parentheses: CV3 jackknife s.e.\ by developer (Obs.; bias: of Obs.\ $-$ Exp., with the benchmark's s.e.\ added in quadrature) and the design-conditional wild-cluster bootstrap s.e.\ of the surface (9,999 draws; Exp.). Brackets: the largest of three few-cluster $p$-values for zero bias that include the benchmark's uncertainty (CV3 with $t(G-1)$; restricted wild-cluster bootstrap-$t$, Webb weights, 9,999 draws, studentized by CR1 or by CV3). Effective numbers of clusters: 1.2--3.8 of 6 in Panel A (no $p$-value is reliable there); 2.5--6.7 of 19 in Panel B. ARC-Challenge and Winogrande: Tobit surface censored at the floor (52 and 16 of 86 runs are at it), level matched to the observed outputs, counterfactual clipped at the same floor. Notability-propensity control (selection on observables): quadratic in the estimated propensity of Epoch notability; Panel A has no notable model.
\end{tablenotes}
\begin{tablenotes}[Source]
Authors' calculations. Observational outputs: Open LLM Leaderboard v1 via \citet{ruan2024observational} and \citet{maiapolo2024sloth}; parameter counts from Hugging Face; release dates from \citet{epochai2026data}. Experimental runs: \citet{bhagia2024establishing}, \citet{teamolmo2024olmo} (Panel B only) and \citet{gadre2024language}.
\end{tablenotes}
\end{table}


\begin{figure}[tp]
\centering
\includegraphics[width=\textwidth]{ra4_obsfix_fig_matched}
\caption{Observational Estimates against the Experimental Technology on the Same Design}\label{fig:app-obs-matched}
\begin{figurenotes}
Bias of each observational estimator relative to the same estimator applied to the outputs of the experimental
technology on the same models (Obs.\ $-$ Exp.), with 95 percent intervals ($\pm1.96$ CV3 standard errors of the bias,
including the benchmark's uncertainty). Blue: strict benchmark (OLMo ladder and \citealp{gadre2024language}; 20 models).
Orange: benchmark including the released OLMo-2 7B and 13B models (57 models). Panels (a) and (b): HellaSwag,
compute-only elasticity $\theta_C$ and data elasticity $\theta_D$. Panel (c): pooled OLS bias in $\theta_C$ by output;
ARC-Challenge and Winogrande use the censored (Tobit) benchmark. Estimates and conservative few-cluster $p$-values are
in Table~\ref{tab:app-matched}.
\end{figurenotes}
\end{figure}

\emph{Returns to compute.}---Once the design is matched, cross-lab returns to compute on HellaSwag do not differ
detectably from the experimental technology (Table~\ref{tab:app-matched}, Figure~\ref{fig:app-obs-matched}). On the
57-model design, pooled OLS gives $\theta_C=0.328$ against 0.329 from the experimental technology on the same models,
a bias of $-0.001$ (0.041) with $p=0.98$. Across six estimators, from pooled OLS to lab-by-period fixed effects and a
notability-propensity control, the biases lie between $-0.029$ and $+0.037$, and every $p$-value is at least 0.36. The
strict surface evaluated on the same 57 models, an extrapolation for 37 of them, gives $+0.014$ (0.045), so the result
does not rest on anchoring the surface with OLMo-2. On the strict support itself the point estimates do not move (from
$-0.024$ to $+0.037$ across five estimators), but with six clusters they carry no inferential weight.

This is a failure to reject in a low-power joint test, not evidence of unbiasedness. On 57 models the 95 percent
interval for the OLS bias is $[-0.088,0.086]$ with $t(18)$ critical values, about $\pm26$ percent of the benchmark. The
bias detectable with 80 percent power is 0.12, or 37 percent of the benchmark; on the strict support it is 0.27, or 62
percent. The test covers only models with at most 8.8B parameters, and its benchmark is the technology of the AI2 and
OpenLM recipes.

\emph{Other benchmarks.}---ARC-Challenge and Winogrande are less informative, because many small designed runs score
at or below chance: 52 of the 86 runs are at the floor of chance-adjusted accuracy on ARC-Challenge and 16 on
Winogrande, against 2 on HellaSwag. A least-squares surface fitted to the clipped log-odds makes the observational
compute elasticity on ARC-Challenge look attenuated, with a bias of $-0.098$ on 57 models (as reported in an earlier
version of this paper). The attenuation does not survive
modeling the censoring. We fit a Tobit surface censored at the floor, match its level to the observed outputs, and clip
the counterfactual at the same floor; the bias is then $+0.036$ (0.069), with $p=0.61$ (Table~\ref{tab:app-matched},
Panel~C). The Tobit benchmark extrapolates a Gaussian latent surface below the floor for 60 percent of the runs, so we
do not read it as showing that the attenuation was an artifact. Winogrande keeps a borderline attenuation: $-0.150$
(0.077) on 57 models, with $p$-values between 0.068 and 0.084 across the three procedures. On the strict support the
Winogrande estimate is similar but uninformative ($-0.179$, $p=0.53$). The three-benchmark composite gives $-0.018$
(0.050).

Two parts of the joint null are not tested. Observational outputs come from the Open LLM Leaderboard (HellaSwag 10-shot,
ARC-Challenge 25-shot, Winogrande 5-shot), whereas the designed runs were evaluated in the OLMES and LLM-foundry
harnesses with other prompts and shot counts; we harmonize them only through the log-odds of chance-adjusted accuracy.
And benchmark contamination, or classifier-filtered corpora that favor benchmark-like text, raise output at given
inputs; family effects absorb such shifts only when they are common to a family.

\emph{The split between parameters and data.}---The split is not reliably estimated. With family fixed effects on the
57-model design, the observational parameter elasticity exceeds the benchmark by 0.203 (0.059) and the data elasticity
falls short of it by 0.262 (0.106). The conservative $p$-values are 0.048 and 0.027, and dropping any one developer
leaves both signs unchanged (Table~\ref{tab:app-inference}). But the identifying variation is thin. In the full panel
the within-family standard deviation of $\ln D$ is 0.42, against 1.29 for $\ln N$. The effective number of clusters is
about 4 of 19, and two families supply 65 percent of the leverage for $\theta_D$: Cerebras-GPT, trained at exactly 20
tokens per parameter, and Phi, trained on synthetic data. Under the strict benchmark the biases have similar size
($+0.185$ and $-0.227$) but conservative $p$-values of 0.14 and 0.15. Biases of opposite sign are also hard to
reconcile with transmission. If wedges are unrelated to productivity at given compute, transmission shifts both
elasticities by the same amount (Section~\ref{app:obs-theory}). Opposite signs point instead to factor-biased recipe
changes within families or to errors in $D$, which the benchmark cannot distinguish from bias. We read them as a
within-family disagreement driven by two families, not as a finding.

\begin{table}[tp]
\centering
\caption{Few-Cluster Inference for the Family Fixed-Effects Biases}
\label{tab:app-inference}
\footnotesize
\setlength{\tabcolsep}{3pt}
\begin{tabular}{@{}l ccc ccc@{}}
\toprule
 & \multicolumn{3}{c}{$\theta_N$ bias $=0.203$} & \multicolumn{3}{c}{$\theta_D$ bias $=-0.262$} \\
\cmidrule(lr){2-4}\cmidrule(lr){5-7}
 & SE Obs. & SE bias & $p$ & SE Obs. & SE bias & $p$ \\
\midrule
CR1, all $k$ (earlier version), normal & 0.082 & 0.086 & 0.019 & 0.161 & 0.163 & 0.11 \\
CR1, singletons dropped, normal & 0.071 & 0.077 & 0.008 & 0.140 & 0.142 & 0.065 \\
CR1, nested $k$, $t(18)$ & 0.055 & 0.063 & 0.005 & 0.110 & 0.112 & 0.031 \\
CV3 jackknife, $t(18)$ & 0.052 & 0.059 & 0.003 & 0.104 & 0.106 & 0.024 \\
\addlinespace[2pt]
WCR, benchmark fixed (earlier version) &  &  & 0.046 &  &  & 0.024 \\
WCR, benchmark fixed, nested $k$ &  &  & 0.044 &  &  & 0.025 \\
WCR, with benchmark draws &  &  & 0.048 &  &  & 0.028 \\
WCU, with benchmark draws &  &  & $<$0.0001 &  &  & 0.010 \\
WCR, family weights &  &  & 0.031 &  &  & 0.024 \\
WCR, model weights &  &  & 0.010 &  &  & 0.030 \\
WCR, CV3-studentized, with benchmark draws &  &  & 0.048 &  &  & 0.011 \\
\addlinespace[2pt]
Effective clusters $G^*$ (of 19) & \multicolumn{3}{c}{4.3} & \multicolumn{3}{c}{4.1} \\
Bias, one developer left out & \multicolumn{3}{c}{$[0.157,0.224]$} & \multicolumn{3}{c}{$[-0.339,-0.193]$} \\
\bottomrule
\end{tabular}

\begin{tablenotes}
Family fixed-effects rows of Table~\ref{tab:app-matched}, Panel B (HellaSwag, 57 models, benchmark including OLMo-2); only the inference differs across rows. `Earlier version': the procedures of an earlier version of this paper. SE (bias) adds the design-conditional bootstrap standard error of the benchmark (0.029 for $\theta_N$, 0.020 for $\theta_D$). The CR1 small-sample factor $\frac{G}{G-1}\frac{n-1}{n-k}$ is 2.24 on the variance when $k=32$ counts the family dummies (15 of the 30 families are singletons), 1.64 with singletons dropped, and 1.04 when fixed effects nested within developer clusters are not counted \citep{cameron2015practitioner}. A bootstrap-$t$ is invariant to this constant; a normal reference is not, which is why the earlier version paired a $p$-value of 0.02 with a $t$-ratio of 1.6 for $\theta_D$. WCR/WCU: restricted/unrestricted wild-cluster bootstrap-$t$, developer clusters, Webb weights, 9,999 draws \citep{cameron2008bootstrap,webb2023reworking}; subcluster: weights drawn by family or by model with developer-clustered standard errors \citep{mackinnon2018wild}; CV3-studentized: the bootstrap statistic uses the CV3 standard error of $\text{Obs.}-\text{Exp.}$ in the sample and in every draw. CV3: jackknife by developer \citep{mackinnon2023cluster}. $G^*$: \citet{carter2017asymptotic} with $\rho=0$. Identifying variation for $\theta_D$ (partial leverage) by developer: Cerebras 37\%, Microsoft 28\%, Meta 12\%, BigScience 7\%; 7 of 19 developers contribute nothing (they have only singleton families) and 2 more under 0.1\%; by family: Cerebras-GPT 37\%, Phi 28\%, OPT 9\%.
\end{tablenotes}
\begin{tablenotes}[Source]
Authors' calculations; see Table~\ref{tab:app-matched}.
\end{tablenotes}
\end{table}


Table~\ref{tab:app-inference} also explains why an earlier version of this paper paired a $t$-ratio of 1.6 with a
wild-bootstrap $p$-value of 0.02 for the same $\theta_D$ bias. The benchmark's uncertainty is not the cause; its
standard error (0.020) is small against the observational one. The cause is the small-sample factor of the CR1
variance. That version counted the 29 family dummies in $k$ although families are nested within developer clusters,
which raises $(n-1)/(n-k)$ from 1.04 to 2.24 \citep{cameron2015practitioner}. The standard error was therefore 1.47
times too large; a bootstrap-$t$ is invariant to that constant, and a normal reference is not. With the nested
convention, the CV3 jackknife or any restricted wild-cluster variant, the $p$-values are 0.01 to 0.03 for $\theta_D$ and
0.003 to 0.05 for $\theta_N$.

\emph{A Hall-type bias.}---Both worlds agree that $\theta_N>\theta_D$ on this design: 0.42 against 0.27 under the
experimental technology on the 57 models, and $\theta_N-\theta_D=0.21$ (0.08) in the observational data on all 128.
Over-trained models have wedges above one, so their output elasticity with respect to parameters exceeds that with
respect to data. A regression on $\ln C=\ln6ND$ alone therefore imposes equal weights that the data reject, and the
residual of such an equal-weight index misattributes $\tfrac12(\theta_N-\theta_D)(\Delta n-\Delta d)$ to productivity.
This is the Hall-type bias of Corollary~\ref{cor:hall}: as models shift further toward tokens, measured productivity
understates progress.

\emph{How large could transmission bias be?}---We treat the 25 pretraining recipes of \citet{magnusson2025datadecide}
as developers with real outcomes and simulate how they choose model size. Pooled OLS is then unbiased under exogenous
budgets. It overstates $\theta_C$ by 0.044 when funding rises with productivity rank, and understates it by 0.18 (40
percent) when each developer trains to a common capability target. Cross-family productivity dispersion in the
observational panel is about five times that of the recipes, which scales the funding bias to about $+0.2$ and the
target bias beyond $-0.18$. The 95 percent interval of the matched OLS bias excludes both. The data are therefore
consistent with budgets roughly orthogonal to productivity, or with offsetting funding and release selection; the
simulated rules are stylized. In the simulated industry of Online Appendix~\ref{app:mc}, pooled estimation overstates
the frontier elasticity by 13 percent under funding, understates it by 30 percent under a capability target, and
overstates it by 66 percent with one flagship per developer and generation. Estimation within developer and generation,
the observational analogue of an IsoFLOP sweep, is unbiased under every rule, but it requires at least two released
siblings.

% -----------------------------------------------------------------------------------------------------
\subsection{Measurement of Compute and the Standard Remedies}\label{app:obs-compute}

\emph{Measurement error.}---Measurement error in compute is small, but showing it requires an independent measure.
Compute in the panel is $C=6ND$ from reported $N$ and $D$, and Epoch's compute figures for these models are mostly
$6ND$ itself (operation counting) or developer-reported counts of the same kind \citep{epochai2026data}. Only
hardware-time records measure compute independently of the reported inputs. We rebuild compute from them as chip-hours
times peak dense 16-bit throughput times a common utilization of 0.3. Epoch's utilization field is not used, because it
is back-calculated from Epoch's own compute estimate in every record we checked. Across the 27 models with such records, from 12 families,
$\ln(C_{\text{hw}}/6ND)$ has standard deviation 0.29, with 95 percent interval $[0.20,0.41]$
(Table~\ref{tab:app-reliability}). This bounds the error in $\ln6ND$ from above, because it also contains differences
in true utilization and errors in chip-hours. The implied reliability of $\ln C$ on the 57-model design is at least
0.98 pooled and 0.955 within families (0.96 and 0.91 at the upper end of the interval). Correcting for errors in
variables raises the pooled compute elasticity by 0.007 and the within-family one by 0.021, or by 0.042 at the upper
end (Table~\ref{tab:app-compute}).

\input{tables/appE_reliability}

An earlier version of this paper reported a within-family errors-in-variables estimate of 0.858 against a benchmark of
0.409. It combined two errors. It used Epoch's `Confident' uncertainty class ($\sigma_u=\ln3/1.645=0.67$), which
describes the uncertainty of Epoch's own estimate rather than the error in $6ND$. And it computed the within-family
variance of $\ln C$ with an $n-1$ denominator, ignoring the 30 family means. With the correct variance and the
hardware-time bound, the within-family bias is $+0.052$ (0.049), with $p=0.30$.

\emph{Reverse regression.}---The reverse regression of $\ln C$ on output gives an inverse slope of 0.560. By
construction this equals the forward slope, 0.328, divided by the forward $R^2$ of 0.585 (with family effects,
$0.441/0.760=0.580$). Its distance from the same statistic computed on the experimental outputs, 0.352, therefore
measures output variation not explained by compute, that is, productivity dispersion and noise, not a bias. Under
funding the forward--reverse pair need not bracket the truth at all (Proposition~\ref{prop:A-transmission}(iii)).

\begin{table}[tp]
\centering
\caption{Compute-Only Estimators: Measurement Error, Reverse Regression and Instruments}
\label{tab:app-compute}
\footnotesize
\setlength{\tabcolsep}{3pt}
\begin{tabular}{@{}l cccc l@{}}
\toprule
 & Obs. & Exp. & Bias & $\lambda$, $R^2$, $F$ & AR set; $\hat\theta/R^2$ \\
 & (1) & (2) & (3) & (4) & (5) \\
\midrule
\multicolumn{6}{@{}l}{\textit{Panel A. HellaSwag, 57 models (benchmark including OLMo-2)}} \\
OLS, EIV, $\sigma_u=0.29$ & 0.335 & 0.329 & 0.006 & 0.980 &  \\
 & (0.038) & (0.007) & (0.039) [0.90] &  &  \\
Family FE, EIV, $\sigma_u=0.29$ & 0.461 & 0.409 & 0.052 & 0.955 &  \\
 & (0.047) & (0.013) & (0.049) [0.30] &  &  \\
Family FE, EIV, $\sigma_u=0.67$ & 0.576 & 0.409 & 0.166 & 0.766 &  \\
 & (0.059) & (0.013) & (0.060) [0.12] &  &  \\
Family FE, EIV, $\sigma_u=0.67$, $n-1$ & 0.858 & 0.409 & 0.448 & 0.514 &  \\
 & (0.088) & (0.013) & (0.089) [0.016] &  &  \\
Reverse regression & 0.560 & 0.352 & 0.209 & 0.585 & $0.328/0.585$ \\
 & (0.126) & (0.006) &  &  &  \\
Reverse, family FE & 0.580 & 0.432 & 0.147 & 0.760 & $0.441/0.760$ \\
 & (0.122) & (0.012) &  &  &  \\
IV, frontier FLOP/\$ ($n=57$) & 0.724 & 0.012 & 0.712 & 0.90 & $(-\infty,\infty)$ \\
 &  & (0.047) &  &  & $(-\infty,\infty)$ \\
IV, frontier, developer FE ($n=57$) & 8.535 & $-$2.287 & 10.822 & 0.01 & $(-\infty,\infty)$ \\
 &  & (0.190) &  &  & $(-\infty,\infty)$ \\
IV, own hardware ($n=10$) & $-$0.090 & 0.195 & $-$0.285 & 14.3 & $[-1.50,0.34]$ \\
 &  & (0.014) &  &  & $[-0.94,0.27]$ \\
\addlinespace[3pt]
\multicolumn{6}{@{}l}{\textit{Panel B. HellaSwag, 20 models (strictly experimental benchmark)}} \\
OLS, EIV, $\sigma_u=0.29$ & 0.417 & 0.430 & $-$0.013 & 0.973 &  \\
 & (0.019) & (0.007) & (0.020) [0.44] &  &  \\
Family FE, EIV, $\sigma_u=0.29$ & 0.488 & 0.464 & 0.024 & 0.969 &  \\
 & (0.055) & (0.010) & (0.056) [0.65] &  &  \\
Reverse regression & 0.551 & 0.442 & 0.109 & 0.737 & $0.406/0.737$ \\
 & (0.129) & (0.006) &  &  &  \\
Reverse, family FE & 0.505 & 0.476 & 0.029 & 0.937 & $0.473/0.937$ \\
 & (0.081) & (0.009) &  &  &  \\
IV, frontier FLOP/\$ ($n=20$) & 0.008 & 0.384 & $-$0.376 & 5.15 & $(-\infty,\infty)$ \\
 &  & (0.014) &  &  & $[-1.85,0.27]$ \\
IV, frontier, developer FE ($n=20$) & $-$0.237 & 0.332 & $-$0.569 & 3.09 & $(-\infty,\infty)$ \\
 &  & (0.017) &  &  & $(-\infty,\infty)$ \\
IV, own hardware & \multicolumn{5}{l}{fewer than 8 models with the instrument} \\
\bottomrule
\end{tabular}

\begin{tablenotes}
Compute-only estimators of $\theta_C$ on the design-matched samples of Table~\ref{tab:app-matched}. EIV: slope divided by the reliability $\lambda=1-\sigma_u^2/\text{Var}(\ln C)$ (column 4), with $\text{Var}(\ln C)$ pooled or within family; the within-family variance is $\sum\tilde c_i^2/(n-F)$, net of the $F$ family means. $\sigma_u=0.29$ is the standard deviation of $\ln(C_{\text{hw}}/6ND)$ across the 27 main-sample models with hardware-time records (Table~\ref{tab:app-reliability}), an upper bound on the error in $\ln 6ND$; $\sigma_u=\ln3/1.645$ is Epoch's `Confident' class, which describes the uncertainty of Epoch's own estimate, not of $6ND$. `$n-1$': the within-family variance computed with an $n-1$ denominator, as in an earlier version of this paper (reliability 0.514 instead of 0.766 for the same $\sigma_u$). Exp.\ applies the same estimator to the experimental outputs, which are built from the same measured inputs and need no correction. Brackets: restricted wild-cluster bootstrap-$t$ $p$-values for zero bias (CR1-studentized, developer clusters, Webb weights, 9,999 draws, each with an independent draw of the benchmark). Reverse regression: $1/b$ from regressing $\ln C$ on the output. By construction $1/b=\hat\theta_C/R^2$ of the forward regression (column 5), so its gap to Exp.\ measures the output variance not explained by compute (productivity dispersion and noise), not a bias. IV: $\ln C$ instrumented by the log FLOP/s per dollar of the best data-center accelerator released six months before the model, or of the model's own training hardware (Epoch); $F$: cluster-robust first-stage statistic; Anderson--Rubin sets by inversion over $\theta_C\in[-3,4]$ with $t(G-1)$ critical values (second line: normal critical values); $(-\infty,\infty)$: the set covers the whole grid.
\end{tablenotes}
\begin{tablenotes}[Source]
Authors' calculations; see Table~\ref{tab:app-matched}; hardware records from \citet{epochai2026data}.
\end{tablenotes}
\end{table}


\emph{Instruments, panels and proxies.}---The standard remedies for transmission bias cannot be used, for reasons that
concern the data rather than the estimators. The price-performance of the best available data-center accelerator is a
step function of calendar time. On the 57-model design its first-stage $F$ is 0.90 (0.01 with developer fixed effects),
and the Anderson--Rubin sets are unbounded. Applied to the experimental outputs, the same instrument recovers 0.012
against a design-matched elasticity of 0.33: the instrument fails, not the data. The price-performance of each model's
own training hardware is strong only on 10 models from 7 developers ($F=14.3$), and it takes three values, so it acts
as a vintage indicator; its Anderson--Rubin set is $[-1.50,0.34]$. An export-control instrument (Chinese developers
after October 2023, with developer and year effects) is strong among the 90 models with at most 14B parameters and
5T tokens ($F=10.6$). It gives $\theta_C=0.90$ (0.25), with Anderson--Rubin set $[0.49,2.23]$, which is 1.7 to 2.3 times the
least-squares and fixed-effects estimates on that sample. Its exclusion restriction fails because recipes and data
changed at the same time. Dynamic-panel estimators \citep{arellano1991some,blundell1998initial} and the moments of
Proposition~\ref{prop:A-transmission}(v) need three consecutive generations, which only six product-line cells have.
Proxy inversion is unavailable for the reasons given in Section~\ref{app:obs-theory}. The notability-propensity control
of Table~\ref{tab:app-matched} (a quadratic in the estimated propensity of Epoch notability) is selection on
observables, not an Olley--Pakes correction; it barely moves the matched bias ($-0.015$, $p=0.80$).

% -----------------------------------------------------------------------------------------------------
\subsection{Productivity Dispersion}\label{app:obs-tfp}

How dispersed is productivity across model families? The answer depends on the units in which it is measured
(Table~\ref{tab:app-tfp-units}). Let $\Delta$ be the gap between the 90th and 10th percentiles of family productivity
in HellaSwag log-odds, net of inputs. In output units, the 90--10 ratio of the odds, $e^{\Delta}$, is 1.9 to 3.0 across
samples and netting methods. In reducible-loss units, $e^{\gamma\Delta/\theta_C}$ with the frontier elasticity
$\gamma=0.178$ of \citet{besiroglu2024chinchilla}, it is 1.3 to 1.8. Both ranges bracket the 90--10 ratio of 1.92
within US manufacturing industries \citep[see also \citealp{syverson2011determines}]{syverson2004product}. In
compute-equivalent units, $e^{\Delta/\theta_C}$, the compute that the 10th-percentile family needs to match the 90th,
the ratio is 7 to 24 with the 57-model benchmark $\theta_C=0.33$, or 4 to 13 with the strict benchmark's 0.43. Within
developers, in the design of \citet{mertens2026secret}, it is 11.1 in compute units (90 percent interval
$[4.1,18.1]$), against their 41 on another benchmark.

Manufacturing's ratio is both an output-unit and an input-unit ratio, because returns to scale are near one. For
language models the two differ by the factor $1/\theta_C$ in logs, between 2.3 and 3, because returns to compute are
small. The comparison with manufacturing is therefore a choice of cardinalization, not a finding.

\input{tables/appE_tfp}

Three caveats apply. First, family effects absorb benchmark contamination and the similarity of training data to the
benchmark, as well as productivity. Second, training compute omits the compute of teachers and synthetic-data
generators. With inputs netted within families, the synthetic-data families Phi and SmolLM rank first and third of 38,
so measured productivity credits the student with an intermediate input, the gross-output problem of
\citet{gandhi2020identification}. Third, the conversion to compute rests on a single cardinal elasticity from one
benchmark; the intervals in Table~\ref{tab:app-tfp-units} draw $\theta_C$ from its bootstrap. Transmission, finally,
would make all of these ratios understate dispersion (Proposition~\ref{prop:A-transmission}(ii)).

% -----------------------------------------------------------------------------------------------------
\subsection{Algorithmic Progress}\label{app:obs-progress}

\citet{ho2024algorithmic} estimate that the compute needed to reach a given perplexity halves every 8.4 months, from 231
model--benchmark observations in 144 papers. Their preferred model is the technology of Section~\ref{sec:framework}
with $E=0$ and factor-augmenting productivity growing linearly in time, $\psi_N=g_Nt$ and $\psi_D=g_Dt$:
\[
\ln\text{ppl}=\exp(\alpha_{0b}-\alpha g_Nt-\alpha\ln N)+\exp(\beta_{0b}-\beta g_Dt-\beta\ln D),
\]
where $b$ indexes three benchmarks and $t$ is time in years. Effective compute grows at $g_C=g_N+g_D$ per year and
doubles every $\tau_C=12\ln2/g_C$ months; Hicks neutrality is $\alpha g_N=\beta g_D$.

\emph{The ridge.}---The rate is weakly identified because the sample is not on the expansion path, a problem of the kind
studied by \citet{diamond1978measurement}. To first order the data pin down how fast loss falls at fixed inputs, an
average of $\alpha g_N$ and $\beta g_D$ weighted by the mean shares of the two terms in loss (0.55 and 0.45). The
doubling time instead weights the two rates by $1/\alpha$ and $1/\beta$. The two combinations coincide only on the
expansion path, along which parameters would grow with compute at elasticity $a=0.37$ under their exponents. Off the
path, rate pairs that fit equally well imply different doubling times.

Figure~\ref{fig:app-obs-ridge} shows the resulting ridge. The bootstrap correlation between $\alpha g_N$ and $\beta g_D$
is $-0.73$ with the authors' code and $-0.82$ at the converged estimates. The profile-likelihood interval for $\tau_C$
is $[4.1,40.5]$ months. The parameter-augmenting share $g_N/g_C$ can lie anywhere in $[-0.44,3.37]$ (95 percent
profile interval), and over that range $\tau_C$ runs from 5.3 to 39.6 months (Table~\ref{tab:app-progress}). These
likelihood ratios assume independent Gaussian errors, although observations are clustered by paper. The profile over
$g_N/g_C$ also has a second local minimum near $-0.28$; a missed global minimum could only widen the interval.

\begin{figure}[tp]
\centering
\includegraphics[width=\textwidth]{appE_ridge}
\caption{The Ridge in Estimates of Algorithmic Progress}\label{fig:app-obs-ridge}
\begin{figurenotes}
Data of \citet{ho2024algorithmic}, 231 model--benchmark observations. Panel (a): 95 percent joint profile-likelihood
region for the year terms $(\alpha g_N,\beta g_D)$, bootstrap draws from the authors' code, the optimizer's iterates
from its default stopping point to the converged minimum of the same objective, the unpenalized least-squares minimum,
and the Hicks-neutral line. Panel (b): doubling time of effective compute $\tau_C$ along the floor of the profile; the
thick segment lies inside the 95 percent region. Panel (c): profile likelihood-ratio statistic over the
parameter-augmenting share $g_N/g_C$, with the implied $\tau_C$ in months at selected points; the dashed line is the 95
percent cutoff and the shaded band the 95 percent interval, $[-0.44,3.37]$. Likelihood ratios assume independent
Gaussian errors.
\end{figurenotes}
\end{figure}

\emph{Neutrality.}---Precision about the rate is bought with an assumption about the bias. Hicks neutrality is not
rejected (bootstrap $p=0.29$ and 0.94 for the converged penalized and the unpenalized estimators), which reflects low
power rather than evidence of neutrality. Imposing it gives 9.2 months with unpenalized least squares, with interval
$[6.3,14.3]$, or 11.6 months with the converged penalized estimator: the tightest intervals among the specifications
that the data do not reject. As \citet{diamond1978measurement} predict, the rate is identified here by the restriction
on the bias. Even that restriction is fragile. If the true technology has $E>0$, neutrality imposed within the $E=0$
form is misspecified. In simulations on Ho et al.'s design with a true doubling time of 12 months, it yields 4.2 to 4.4
months, so the neutral estimates may overstate progress.

\input{tables/appE_progress}

\emph{The published estimate.}---The published point also depends on the optimizer. We reproduce Ho et al.'s estimates
exactly: the same objective value (0.0518), every parameter within $4\times10^{-7}$, and a bootstrap median of 8.44
months under their protocol. The published 8.4 months is that median; the point estimate is 8.7. It is not the
minimizer of the authors' objective, mean squared error plus an $L_1$ penalty. SciPy's default stopping rule halts the
optimizer after 18 iterations, part of the way along the ridge. Run to convergence, the objective falls to 0.0507, all
300 random starts end below the published value, and $\tau_C=6.1$ months with a paper-cluster interval of $[3.0,22.7]$,
about twice as wide as published. Mean squared error changes by only 0.27 percent: almost all of the improvement in
the objective comes from the penalty on constants whose values depend on an arbitrary normalization. Unpenalized least
squares gives 10.2 months, with interval $[4.1,27.6]$. Estimators of the same model thus give 6.1 to 10.2 months, and
the data cannot arbitrate between them. The finding concerns numerical optimization, not the order of magnitude of
progress. It parallels the finding of \citet{besiroglu2024chinchilla} for the third approach of
\citet{hoffmann2022training}, and convergence problems in demand estimation
\citep{knittel2014estimation,dube2012improving}.

\emph{Attenuated exponents.}---Ho et al.'s data exponent is 0.04, against about 0.37 in experiments, and
\citet{whitfill2025note} reads the gap as evidence that progress is overstated about ninefold: dividing their year
coefficients by experimental exponents gives 76.5 months. Imposing experimental exponents in estimation is rejected
(mean squared error rises by a factor of 1.3 to 4.8 across 14 experimental technologies) and implies negative or
implausibly slow progress. Two specification choices explain much of the exponent gap while moving $\tau_C$ much less
(Table~\ref{tab:app-obs-attenuation}). First, fitting the $E=0$ form to the Chinchilla sweep itself lowers the frontier
elasticity from 0.178 to 0.053, close to the sweep's mean elasticity of total rather than reducible loss (0.054). The
like-for-like gap is therefore about 2 to 2.5-fold, not 7 to 9-fold. Dropping $E$ rescales exponents and year
coefficients by similar factors, so rates barely move: imposing a common $E$ from 0 to 1.5 nats per word doubles
$\hat\gamma$ but moves $\tau_C$ only from 10.2 to 12.4 months. Ho et al.'s data themselves prefer $E=0$ within this
form, so this argument rests on the sweep and on simulations; in the simulations the authors' procedure can also stop
at a local optimum that overstates progress. Second, their $D$ is dataset size, while small models were trained for a
median of 62 to 140 epochs. Measuring effective data under repetition \citep{muennighoff2023scaling} raises $\tau_C$ by
a factor of 1.5 to 4, with an unbounded upper confidence limit.

\input{tables/appE_attenuation}

\emph{Scale dependence and frontier regressions.}---Two further questions need designed data. Whether the
compute-equivalent gain of transformers over earlier architectures depends on scale turns on whether their frontier
elasticities differ (Corollary~\ref{cor:ceg}). In Ho et al.'s data the difference is 0.008, with interval
$[-0.012,0.034]$, and bootstrap intervals for the implied gains span six orders of magnitude
(Table~\ref{tab:app-obs-ceg-sahal}). And regressions of loss on compute that omit time effects overstate the compute
elasticity by a factor of 1.15 on all 231 rows and 1.44 among 36 record-setting releases (Corollary~\ref{cor:sahal}).
Among record setters, whose compute moves closely with time, independently measured growth rates imply an algorithmic
share of 0.37 and a predicted factor of 1.58, consistent with the estimate. On all rows the prediction fails, because
compute and time are only moderately correlated (0.62) and the ordinary omitted-variable formula applies instead.

\begin{table}[tp]
\centering
\caption{Compute-Equivalent Gains of Transformers and the Bias of Frontier Regressions}
\label{tab:app-obs-ceg-sahal}
\footnotesize
\setlength{\tabcolsep}{2.5pt}
\begin{tabular}{@{}lcccccc@{}}
\toprule
\multicolumn{7}{@{}l}{\textit{Panel A. Gain of transformers over other architectures}} \\[1pt]
 & \multicolumn{2}{c}{Frontier elasticity} & Difference & \multicolumn{3}{c}{Gain $f(C)$ at $C=$} \\
\cmidrule(lr){2-3}\cmidrule(lr){5-7}
 & Other & Transformer & [95\% CI] & $10^{17}$ & $10^{19}$ & $10^{23}$ \\
\midrule
Penalized & 0.0133 & 0.0217 & 0.0085 [$-$0.012, 0.034] & 0.02 & 0.45 & 163 \\
Unpenalized & 0.0197 & 0.0210 & 0.0013 [$-$0.011, 0.028] & 10.6 & 14.5 & 27.1 \\
\midrule
\multicolumn{7}{@{}l}{\textit{Panel B. Log loss on log compute, with and without a time trend}} \\[1pt]
 & & \multicolumn{2}{c}{Compute slope} & Inflation & Compute & Implied \\
\cmidrule(lr){3-4}
 & Rows & No trend & Trend & [95\% CI] & growth & $s_A$ \\
\midrule
All rows & 231 & 0.0282 (0.0020) & 0.0246 (0.0018) & 1.15 [1.10, 1.21] & 1.12 & 0.13 \\
Record setters & 36 & 0.0329 (0.0024) & 0.0228 (0.0027) & 1.44 [1.24, 1.78] & 2.64 & 0.31 \\
\bottomrule
\end{tabular}
\par\smallskip
\begin{tablenotes}Panel A: model 13 of \citet{ho2024algorithmic} (architecture-specific exponents, common
factor-augmenting year terms, $E=0$), estimated with the penalized objective run to convergence or by unpenalized
NLS, on their 231-row sample, WikiText-103 in 2023. $f(C)$ is the compute a
non-transformer lab needs to match the transformer frontier loss at compute $C$ (FLOP), relative to $C$; with $E=0$,
$\ln f$ is linear in $\ln C$ and constant only if the two frontier elasticities coincide, consistent with the condition
in Corollary~\ref{cor:ceg} (equal $E$ and equal $\gamma$). Non-transformer models span $C\leq3\times10^{18}$ FLOP,
so $f$ above that is extrapolated; bootstrap intervals of $f$ span more than six orders of magnitude. Panel B: OLS of
the log of log perplexity on $\ln(6ND)$ and benchmark indicators, with and without a linear publication-year trend;
paper-clustered standard errors in parentheses; inflation is the ratio of the two slopes, with a paper-cluster
bootstrap interval; compute growth is the within-sample growth of log compute per year; the implied algorithmic share
$s_A$, one minus the ratio of the slope with trend to the slope without, is backed out of the inflation factor and is not an
independent test. Record setters are releases with the best perplexity on the benchmark to date.\end{tablenotes}
\begin{tablenotes}[Source]Data of \citet{ho2024algorithmic}; authors' estimates.\end{tablenotes}
\end{table}


\emph{The missing data.}---Proposition~\ref{prop:A-dmr} names the data that would settle these questions: IsoFLOP
sweeps at several vintages on a common output, whose drift in the compute-optimal $M^*$ signs the bias of technical
change and whose curvature fixes its size. Our controlled experiment (Section~\ref{sec:tech}) supplies the
cross-sectional analogue for a change in data quality at one date.
\textcolor{red}{[TBD-m9: one sentence on the shift in the compute-optimal $M^*$ at fixed compute between the two
corpora of our experiment, and hence the sign of the Hicks bias of the data-quality difference
(Proposition~\ref{prop:A-dmr}(i) applied across corpora rather than dates); delete if m9 cannot answer it.]}

% -----------------------------------------------------------------------------------------------------
\subsection{Allocative versus Technical Change}\label{app:obs-allocative}

Measured progress, the fall in the compute needed to reach a given loss, mixes shifts of the frontier with movements
toward it. By Proposition~\ref{prop:farrell}, the move from the allocations of \citet{kaplan2020scaling} to those of
\citet{hoffmann2022training} changed the mix, not the technology, so its gain is allocative; under the
\citet{besiroglu2024chinchilla} parameters, Gopher's allocation used about twice the compute its loss required.
\citet{gundlach2025origin} define the compute-equivalent gain from such rebalancing and count it as algorithmic
progress.

\emph{Measurement.}---Measuring the gain actually realized between eras requires a choice. Training-cost efficiency, the
Farrell efficiency against the training-only frontier, $C_{\min}(L(N,D))/6ND$, counts deliberate over-training as
waste, although over-training reflects the value that developers place on compact models (Section~\ref{sec:wedge}).
Serving demand cannot produce a wedge below one (Proposition~\ref{prop:wedge}). We therefore count only the
elimination of under-training: training-cost efficiency evaluated at $\min(w,1)$, which within the Chinchilla family
depends on the wedge alone. The truncated measure attributes every wedge below one to allocative error. A wedge below
one can also reveal a binding data constraint or data-augmenting productivity specific to the developer
(Proposition~\ref{prop:A-wedge}(ii) and (iv)). For 2020--21 models trained on a few hundred billion tokens, a belief in the
\citet{kaplan2020scaling} law is the more plausible reading, but it remains an assumption.

\emph{Realized gains.}---We compare models in the Epoch database with at least $10^{23}$ FLOP: 7 released in 2020--21
(era 1) and 92 in 2022--24 (era 2), all dense, trained for at most four epochs, with $D=C/6N$
(Table~\ref{tab:app-allocative}, Figure~\ref{fig:app-obs-allocative}). Every era-1 model has $w<1$ under every
technology, and the median ratio of tokens to parameters rose from 1.7 to 92 between the eras. Counting only the
elimination of under-training, geometric-mean training-cost efficiency rose by a factor of 2.0 under the reference
technology (Chinchilla with $\kappa$ free) and 1.7 under \citet{besiroglu2024chinchilla}, with 95 percent interval
$[1.21,2.44]$. Across all 14 estimated technologies the factor ranges from 1.1 to 4.9. Era-2 releases closed 80 to 94
percent of the era-1 under-training gap, in logs. The positive sign is close to mechanical, because every era-1 model is
under-trained and truncation can only raise era-2 efficiency; the content is in the size. Without truncation the gain
ranges from 0.55 to 4.3 across technologies: counting over-training as waste can reverse its sign.

\input{tables/appE_allocative}

\begin{figure}[tp]
\centering
\includegraphics[width=0.85\textwidth]{ra4_obsfix_fig_allocative}
\caption{Realized Allocative Gains from Rebalancing, by Technology}\label{fig:app-obs-allocative}
\begin{figurenotes}
Ratio of geometric-mean training-cost efficiency, 2022--24 against 2020--21, for Epoch models with at least $10^{23}$
FLOP (7 and 92 models), under each technology, with 95 percent percentile-bootstrap intervals over models within era
(log scale). Orange: efficiency against the training-only frontier, which counts over-training as waste. Blue: the
wedge truncated at one, which counts only the elimination of under-training. Only the \citet{besiroglu2024chinchilla}
interval also reflects uncertainty in the technology's parameters. Estimates in Table~\ref{tab:app-allocative}.
\end{figurenotes}
\end{figure}

The intervals resample models within era; only the Besiroglu interval also draws the technology's parameters, and with
seven era-1 models all intervals are rough. Technologies estimated on other sweeps use their own conventions for $N$,
$D$ and the loss corpus, so their levels are indicative. The Marin fits, whose parameter exponents are about twice
Chinchilla's, give the largest gains (3.5 to 4.9). Samples of the five largest models per year and of Epoch's frontier
models give smaller and noisier truncated gains (1.3 and 1.5 under Besiroglu).

\emph{Comparison with progress.}---Against the effective-compute gain implied by Ho et al.'s published rate over the
same interval, 8.9-fold over 2.28 years, the realized allocative gain amounts to 6 to 73 percent in logs across
technologies (31 percent under the reference technology). The denominator itself ranges from 6.4 to 22.5 across the
estimators of Section~\ref{app:obs-progress}. \citet{gundlach2025origin} report that the Kaplan-to-Chinchilla
rebalancing is worth about a factor of 10 at the 2025 frontier. Our counterfactual agrees: under the Besiroglu
technology, continuing Kaplan's rule from GPT-3 would cost a factor of 1.6 at GPT-3's compute, 5.2 at
$3.8\times10^{25}$ FLOP and 11.7 at $5\times10^{26}$ FLOP. Because Kaplan-rule allocations have $w<1$ under every
technology, truncation does not affect this counterfactual. What we add is the contrast with the gain realized between
eras, 1.1 to 4.9, which is smaller than both. Under Ho et al.'s own exponents, finally, every model has $w>1$. An
attenuated data exponent makes tokens nearly worthless, so a regression with those exponents cannot credit rebalancing
to inputs and loads it onto the time trend.

\emph{Summary.}---Cross-lab data thus deliver returns to compute once the design is matched, though only as a
low-power failure to reject, and measures of productivity dispersion whose size depends on the units. They do not
deliver the split between inputs, the bias of technical change, or its rate without a restriction on the bias. In each
case the missing variation is the transverse variation that designed IsoFLOP sweeps supply.
